\chapter{The de Rham Complex and Discrete Exactness}
\label{ch:derham}

\chapterorient{The previous chapter introduced the four finite-element spaces one
at a time, by the degrees of freedom that define them, and ended with a promise:
that $\Hone$, $\Hcurl$, $\Hdiv$, and $\Ltwo$ are not four unrelated spaces but the
four rungs of a single ladder, linked by the gradient, curl, and divergence. This
chapter keeps that promise. It is the structural heart of Part~II. We assemble the
four spaces into the \emph{de Rham complex}, prove the two vector-calculus
identities $\Curl\Grad=0$ and $\Div\Curl=0$ that make it a complex, and read off
their physical content---that a gradient field is curl-free and a curl field is
divergence-free, the very facts that license the scalar and vector potentials of
\cref{ch:maxwell}. Then comes the payoff that makes the whole apparatus matter for
a simulator: the discrete finite-element spaces inherit the \emph{same} exactness,
through a commuting diagram relating interpolation to differentiation. That single
property---discrete exactness---is why edge elements avoid spurious modes, why the
eigenmode solver can project onto a divergence-free subspace, and why the gradient
fields form the null space the curl--curl operator must learn to ignore. By the end
you will be able to draw the complex from memory, verify exactness by hand on a
small mesh, and recognize the de Rham structure wherever it surfaces in the
solver chapters ahead.}

\section{Four spaces, three arrows}

\Cref{ch:functionspaces} left us with \cref{tab:fourdofs}: four spaces, each
defined by where its degrees of freedom live and what continuity that forces.
Reading the table top to bottom traced a ladder---vertices, edges, faces, cell
interiors; full continuity, tangential, normal, none. We now make the rungs of
that ladder into a single mathematical object by naming the \emph{arrows} between
them.

The three arrows are the differential operators of vector calculus. The gradient
$\Grad$ turns a scalar field into a vector field; the curl $\Curl$ turns a vector
field into another vector field; the divergence $\Div$ turns a vector field into a
scalar field. What \cref{ch:functionspaces} hinted at, and we now assert
precisely, is that each of these operators maps one of our four spaces
\emph{exactly into the next}:
\[
  \Hone \xrightarrow{\;\Grad\;} \Hcurl \xrightarrow{\;\Curl\;} \Hdiv
         \xrightarrow{\;\Div\;} \Ltwo .
\]
This sequence is the \emph{de Rham complex}. It is the central diagram of this
book, and we draw it properly in \cref{fig:derham}.

\begin{figure}[t]
\centering
\begin{tikzpicture}[node distance=24mm and 30mm]
  % --- the four space nodes ---
  \node[space, minimum width=18mm] (h1)   {$\Hone$};
  \node[space, minimum width=18mm, right=of h1]   (hcurl) {$\Hcurl$};
  \node[space, minimum width=18mm, right=of hcurl] (hdiv)  {$\Hdiv$};
  \node[space, minimum width=18mm, right=of hdiv]  (l2)    {$\Ltwo$};
  % --- the three differential arrows ---
  \draw[flow] (h1)    -- node[above,font=\small]{$\Grad$}
                         node[below,font=\scriptsize\itshape,simGray]{gradient} (hcurl);
  \draw[flow] (hcurl) -- node[above,font=\small]{$\Curl$}
                         node[below,font=\scriptsize\itshape,simGray]{curl} (hdiv);
  \draw[flow] (hdiv)  -- node[above,font=\small]{$\Div$}
                         node[below,font=\scriptsize\itshape,simGray]{divergence} (l2);
  % --- DoF annotations under each node ---
  \node[font=\scriptsize, simBlue,   below=10mm of h1.south,    align=center] (a1)
    {nodal\\(vertices)\\value $V$};
  \node[font=\scriptsize, simRust,   below=10mm of hcurl.south, align=center] (a2)
    {edge\\(circulations)\\$\vE,\vA$};
  \node[font=\scriptsize, simGreen,  below=10mm of hdiv.south,  align=center] (a3)
    {face\\(fluxes)\\$\vB,\vD$};
  \node[font=\scriptsize, simViolet, below=10mm of l2.south,    align=center] (a4)
    {cell\\(densities)\\$\rho$};
  \draw[refedge] (h1.south)    -- (a1.north);
  \draw[refedge] (hcurl.south) -- (a2.north);
  \draw[refedge] (hdiv.south)  -- (a3.north);
  \draw[refedge] (l2.south)    -- (a4.north);
  % --- the two exactness identities, above the arrows ---
  \node[font=\scriptsize\itshape, simInk, above=12mm of hcurl.north] (id1)
    {$\Curl\Grad=0$};
  \node[font=\scriptsize\itshape, simInk, above=12mm of hdiv.north]  (id2)
    {$\Div\Curl=0$};
  \draw[refedge, bend right=18] (h1.north)  to (id1);
  \draw[refedge, bend left=18]  (hdiv.north) to (id1);
  \draw[refedge, bend right=18] (hcurl.north) to (id2);
  \draw[refedge, bend left=18]  (l2.north)    to (id2);
\end{tikzpicture}
\caption[The de Rham complex]{The de Rham complex: the four finite-element spaces
of \cref{ch:functionspaces}, linked by the three differential operators. Each space
carries a different \emph{kind} of degree of freedom (blue nodal values, rust edge
circulations, green face fluxes, violet cell densities) and each arrow is a
\emph{first-order} differential operator that maps each space exactly into the next.
The two identities above the arrows---$\Curl\Grad=0$ and $\Div\Curl=0$---are what make
this sequence a \emph{complex}: applying two consecutive arrows annihilates everything.
This single picture organizes the whole of electromagnetic finite-element theory;
it recurs throughout the book.}
\label{fig:derham}
\end{figure}

Before we can call it a complex we must check that the arrows even \emph{type-check}:
that each operator genuinely lands in the next space. The four spaces were defined
in \cref{ch:functionspaces} precisely so that they do.

\begin{itemize}[nosep]
\item \textbf{$\Grad : \Hone \to \Hcurl$.} A field $u\in\Hone$ is square-integrable
  with a square-integrable gradient---that is the very definition of $\Hone$. So
  $\Grad u$ is square-integrable. Is it in $\Hcurl$? A field is in $\Hcurl$ if its
  curl is square-integrable; and $\Curl\Grad u = 0$ (proven below), which is
  certainly square-integrable. So $\Grad u\in\Hcurl$. The tangential trace of
  $\Grad u$ is the tangential derivative of the continuous scalar $u$---continuous
  across faces---so $\Grad u$ has exactly the tangential continuity that $\Hcurl$
  requires.
\item \textbf{$\Curl : \Hcurl \to \Hdiv$.} A field $\vu\in\Hcurl$ has
  square-integrable curl by definition, so $\Curl\vu$ is square-integrable. Is it in
  $\Hdiv$? A field is in $\Hdiv$ if its divergence is square-integrable; and
  $\Div\Curl\vu=0$, square-integrable. So $\Curl\vu\in\Hdiv$. The normal trace of a
  curl is determined by the tangential data of $\vu$ on the face, which $\Hcurl$
  makes continuous---so $\Curl\vu$ has the normal continuity $\Hdiv$ requires.
\item \textbf{$\Div : \Hdiv \to \Ltwo$.} A field $\vu\in\Hdiv$ has square-integrable
  divergence by definition. So $\Div\vu\in\Ltwo$, no constraint to check: $\Ltwo$
  asks for nothing beyond square-integrability.
\end{itemize}

The arrows fit. Each space is built with just enough regularity that the
differential operator leaving it lands in the next space with just enough regularity
to leave again. This is no accident of our four choices: it is the reason these four
spaces, and not others, are the canonical ones. They are the spaces \emph{on which
the de Rham complex closes}.

\begin{notationbox}
We write the complex left to right, $\Hone\to\Hcurl\to\Hdiv\to\Ltwo$, with the
operators $\Grad,\Curl,\Div$ over the arrows. In three dimensions this is the full
complex; in two dimensions it shortens to
$\Hone\xrightarrow{\Grad}\Hcurl\xrightarrow{\operatorname{curl}}\Ltwo$, where the
\emph{scalar} curl $\operatorname{curl}(u_1,u_2)=\partial_x u_2-\partial_y u_1$ plays
the role of both curl and divergence at once. Palace builds the discrete arrows for
both cases; the 2-D scalar curl is the \code{H\_CURL}$\to$\code{INTEGRAL} branch of
its interpolator (\cref{sec:discrete-arrows}). Throughout this chapter ``the
complex'' means the 3-D version unless stated otherwise.
\end{notationbox}

\section{Exactness: range equals kernel}

A sequence of spaces and arrows is a \emph{complex} when the composition of any two
consecutive arrows is zero. It is \emph{exact} when something stronger holds: the
\emph{range} of each arrow is exactly the \emph{kernel} of the next. We take these
two properties in turn, because they are different in strength and different in
consequence, and the discrete theory will need to track them separately.

\subsection{The complex property: two consecutive arrows vanish}

The complex property is the pair of vector-calculus identities
\[
  \Curl\,\Grad = 0, \qquad \Div\,\Curl = 0.
\]
These say that the gradient of any scalar has zero curl, and the curl of any vector
has zero divergence. They are usually quoted as facts; we derive them, because the
derivation is short, it is the engine of everything that follows, and---crucially---
the \emph{discrete} versions will be derived the same way.

\begin{lemma}[The complex identities]\label{lem:complex}
For any twice-continuously-differentiable scalar field $\phi$ and vector field $\vF$,
\[
  \Curl(\Grad\phi) = \mathbf{0}, \qquad \Div(\Curl\vF) = 0.
\]
\end{lemma}

\begin{proof}
Both follow from a single fact: \emph{mixed partial derivatives commute},
$\partial_i\partial_j\phi=\partial_j\partial_i\phi$ (Clairaut's theorem, valid for
$C^2$ fields).

For the first identity, compute the curl of a gradient component by component. With
$\Grad\phi=(\partial_x\phi,\partial_y\phi,\partial_z\phi)$, the $x$-component of the
curl is
\[
  (\Curl\Grad\phi)_x
  = \partial_y(\Grad\phi)_z - \partial_z(\Grad\phi)_y
  = \partial_y\partial_z\phi - \partial_z\partial_y\phi
  = 0,
\]
the last step by commuting the mixed partials. The $y$- and $z$-components vanish
identically the same way. Hence $\Curl\Grad\phi=\mathbf 0$.

For the second identity, with $\vF=(F_1,F_2,F_3)$,
\[
  \Curl\vF
  = \bigl(\partial_y F_3 - \partial_z F_2,\;
          \partial_z F_1 - \partial_x F_3,\;
          \partial_x F_2 - \partial_y F_1\bigr),
\]
and its divergence is
\[
  \Div\Curl\vF
  = \partial_x(\partial_y F_3 - \partial_z F_2)
  + \partial_y(\partial_z F_1 - \partial_x F_3)
  + \partial_z(\partial_x F_2 - \partial_y F_1).
\]
The six terms cancel in pairs: $\partial_x\partial_y F_3$ against
$-\partial_y\partial_x F_3$, $-\partial_x\partial_z F_2$ against
$\partial_z\partial_x F_2$, and $\partial_y\partial_z F_1$ against
$-\partial_z\partial_y F_1$, each pair by commuting the mixed partials. The total is
zero.
\end{proof}

Both identities are corollaries of one symmetry---that the order of differentiation
does not matter. We will revisit this in \cref{wex:component}, where we watch the
cancellation happen on a concrete polynomial field, and again in
\cref{sec:discrete-exact}, where the \emph{discrete} gradient and curl satisfy the
same identities for a discrete reason that is structurally identical.

\subsection{Exactness: the kernel is filled by the previous range}

The complex property says $\operatorname{range}(\Grad)\subseteq\ker(\Curl)$: every
gradient is curl-free, so the range of $\Grad$ sits inside the kernel of $\Curl$.
\emph{Exactness} asserts the reverse inclusion as well---that the kernel contains
\emph{nothing but} previous-range elements:
\[
  \operatorname{range}(\Grad) = \ker(\Curl), \qquad
  \operatorname{range}(\Curl) = \ker(\Div).
\]
In words: \emph{every} curl-free field is a gradient, and \emph{every}
divergence-free field is a curl. \Cref{fig:exact} draws the two pictures side by
side---the easy inclusion (a complex) and the full equality (exactness).

\begin{figure}[t]
\centering
\begin{tikzpicture}[scale=1.0, every node/.style={font=\scriptsize}]
% --- Left: complex (one inclusion) ---
\begin{scope}
  \node[font=\footnotesize\bfseries] at (1.5,2.55) {complex: range $\subseteq$ kernel};
  % kernel of Curl as a big ellipse inside Hcurl
  \draw[simInk,thick] (1.5,1.0) ellipse (1.6 and 1.25);
  \node[simInk] at (1.5,2.05) {$\ker(\Curl)$};
  % range of Grad as a smaller blob inside
  \fill[simBlue!18] (1.2,0.85) ellipse (0.85 and 0.6);
  \draw[simBlue,thick] (1.2,0.85) ellipse (0.85 and 0.6);
  \node[simBlue] at (1.2,0.85) {$\operatorname{range}(\Grad)$};
  \node[simRust,font=\scriptsize] at (2.35,0.25) {gap?};
  \node[simGray,align=center] at (1.5,-0.7) {curl-free fields not\\[-1pt]known to be gradients};
\end{scope}
% --- Right: exactness (equality) ---
\begin{scope}[shift={(5.6,0)}]
  \node[font=\footnotesize\bfseries] at (1.5,2.55) {exact: range $=$ kernel};
  \draw[simInk,thick] (1.5,1.0) ellipse (1.6 and 1.25);
  \node[simInk] at (1.5,2.05) {$\ker(\Curl)$};
  \fill[simBlue!18] (1.5,1.0) ellipse (1.55 and 1.2);
  \draw[simBlue,thick] (1.5,1.0) ellipse (1.55 and 1.2);
  \node[simBlue] at (1.5,1.0) {$\operatorname{range}(\Grad)$};
  \node[simGreen,font=\scriptsize] at (1.5,-0.05) {no gap};
  \node[simGray,align=center] at (1.5,-0.7) {every curl-free field\\[-1pt]\emph{is} a gradient};
\end{scope}
\end{tikzpicture}
\caption[Complex versus exact]{The difference between a complex and an exact
sequence, drawn at the $\Hcurl$ rung. \emph{Left:} the complex property
$\Curl\Grad=0$ guarantees only that the range of $\Grad$ (blue) sits \emph{inside}
the kernel of $\Curl$ (the curl-free fields); a gap may remain---curl-free fields
that are not gradients. \emph{Right:} \emph{exactness} closes the gap: the range
fills the kernel exactly, so every curl-free field is a gradient. The size of the
gap is a property of the domain's \emph{topology} (\cref{rmk:cohomology}); for the
simply-connected domains of most simulations the gap is empty and the sequence is
exact.}
\label{fig:exact}
\end{figure}

The reverse inclusions are not automatic---and whether they hold is a statement
about the \emph{shape} of the domain, not about calculus. On a domain with no holes
(a ball, a box, any \emph{simply-connected} region with connected boundary) the
sequence is exact: a curl-free field on such a domain really is a gradient, and a
divergence-free field really is a curl. On a domain with a hole---a solid torus, the
region around a wire---the reverse inclusion can fail, and the failure is measured by
the topology. That refinement is the subject of \cref{rmk:cohomology}; for now we
record the clean statement for the domains where it holds.

\begin{theorem}[Exactness on a simply-connected domain]\label{thm:exact}
Let $\dom\subset\Real^3$ be a bounded, simply-connected domain with connected
boundary. Then the de Rham complex on $\dom$ is exact:
\[
  \operatorname{range}(\Grad)=\ker(\Curl), \qquad
  \operatorname{range}(\Curl)=\ker(\Div).
\]
Equivalently: every curl-free field is the gradient of a scalar potential, and every
divergence-free field is the curl of a vector potential.
\end{theorem}

The proof constructs the potential explicitly---one integrates the field along paths
from a base point, and simple-connectivity guarantees the answer is
path-independent. We do not reproduce it (it is the Poincar\'e-lemma computation of
any vector-calculus course); the references in \cref{sec:furtherreading} give it in
full. What matters here is not the proof but the \emph{use}: exactness is the
mathematical fact behind the electromagnetic potentials, and it is the property the
discrete spaces must preserve.

\section{The potentials of Maxwell are de Rham facts}
\label{sec:potentials}

\Cref{ch:maxwell} introduced two potentials almost in passing: the scalar potential
$V$ with $\vE=-\Grad V$, and the magnetic vector potential $\vA$ with $\vB=\Curl\vA$.
We can now see that neither was a lucky guess. Each is forced---\emph{exists at
all}---by the exactness of the de Rham complex.

Two of Maxwell's equations are homogeneous (source-free) in the regimes the
simulator solves:
\[
  \Curl\vE = 0 \quad\text{(electrostatics)}, \qquad
  \Div\vB = 0 \quad\text{(always).}
\]
Read each through \cref{thm:exact}. The electrostatic field $\vE$ is curl-free; by
exactness at the $\Hcurl$ rung, a curl-free field is a gradient, so there exists a
scalar $V$ with $\vE=-\Grad V$ (the sign is convention). The magnetic field $\vB$ is
divergence-free---always, this is the ``no magnetic monopoles'' law; by exactness at
the $\Hdiv$ rung, a divergence-free field is a curl, so there exists a vector field
$\vA$ with $\vB=\Curl\vA$. \Cref{fig:potentials} lays the correspondence out.

\begin{figure}[t]
\centering
\begin{tikzpicture}[scale=1.0, every node/.style={font=\footnotesize}]
% --- top row: electrostatic / scalar potential ---
\begin{scope}
  \node[stage, minimum width=30mm] (e1) {$\Curl\vE=0$};
  \node[font=\scriptsize\itshape,simGray, right=8mm of e1] (e2) {exactness at $\Hcurl$};
  \node[stage, minimum width=30mm, right=8mm of e2] (e3) {$\vE=-\Grad V$};
  \draw[flow] (e1) -- (e2);
  \draw[flow] (e2) -- (e3);
  \node[font=\scriptsize,simBlue, below=1mm of e1] {$\vE$ curl-free};
  \node[font=\scriptsize,simBlue, below=1mm of e3] {scalar potential $V\in\Hone$};
\end{scope}
% --- bottom row: magnetostatic / vector potential ---
\begin{scope}[shift={(0,-2.4)}]
  \node[stage, minimum width=30mm] (b1) {$\Div\vB=0$};
  \node[font=\scriptsize\itshape,simGray, right=8mm of b1] (b2) {exactness at $\Hdiv$};
  \node[stage, minimum width=30mm, right=8mm of b2] (b3) {$\vB=\Curl\vA$};
  \draw[flow] (b1) -- (b2);
  \draw[flow] (b2) -- (b3);
  \node[font=\scriptsize,simGreen, below=1mm of b1] {$\vB$ divergence-free};
  \node[font=\scriptsize,simGreen, below=1mm of b3] {vector potential $\vA\in\Hcurl$};
\end{scope}
\end{tikzpicture}
\caption[Potentials are exactness]{The electromagnetic potentials of
\cref{ch:maxwell} read as de Rham facts. \emph{Top:} the electrostatic Maxwell
equation $\Curl\vE=0$ says $\vE$ is curl-free; exactness at the $\Hcurl$ rung
(\cref{thm:exact}) supplies a scalar potential $V$ with $\vE=-\Grad V$, living one
rung up the complex in $\Hone$. \emph{Bottom:} the law $\Div\vB=0$ says $\vB$ is
divergence-free; exactness at the $\Hdiv$ rung supplies a vector potential $\vA$ with
$\vB=\Curl\vA$, one rung up in $\Hcurl$. The potential always lives in the space
\emph{one rung above} the field it generates---a structural fact, not a coincidence.}
\label{fig:potentials}
\end{figure}

\begin{keyidea}
The potentials are not tricks; they are exactness. ``$\vE$ is curl-free, therefore
$\vE=-\Grad V$'' \emph{is} the statement $\ker(\Curl)=\operatorname{range}(\Grad)$,
read at a particular field. ``$\vB$ is divergence-free, therefore $\vB=\Curl\vA$''
\emph{is} $\ker(\Div)=\operatorname{range}(\Curl)$. The potential of a field always
lives one rung \emph{up} the de Rham complex: $V\in\Hone$ generates $\vE\in\Hcurl$;
$\vA\in\Hcurl$ generates $\vB\in\Hdiv$. When the magnetostatic analysis of
\cref{ch:targets} solves the curl--curl equation $\Curl(\nu\Curl\vA)=\vJ$ for $\vA$
in $\Hcurl$, it is solving for a potential whose \emph{existence} the de Rham complex
guarantees and whose \emph{home space} the complex names.
\end{keyidea}

There is a consequence we will need repeatedly. Because $\vA=\Grad\psi$ added to any
$\vA$ leaves $\Curl\vA$ unchanged (since $\Curl\Grad\psi=0$), the vector potential is
not unique: $\vA$ and $\vA+\Grad\psi$ describe the same $\vB$ for any scalar $\psi$.
This freedom is the \emph{gauge}, and the gradient fields $\Grad\psi$ are precisely
the directions in which it acts. The same gradient fields are the null space of the
curl--curl operator $\Curl(\nu\Curl\,\cdot)$---for $\Curl(\nu\Curl(\Grad\psi))=0$---
which is the reason that operator is singular and the reason the eigenmode solver
must project them out. We return to this in \cref{sec:divfree}.

\section{The discrete complex}
\label{sec:discrete-arrows}

Everything so far has been about the continuous spaces. A simulator computes in the
\emph{discrete} subspaces $\Hone_h,\Hcurl_h,\Hdiv_h,\Ltwo_h$ of
\cref{ch:functionspaces}---the finite-dimensional spaces of nodal, edge, face, and
cell functions. The question that makes or breaks the method is whether the discrete
spaces form a complex of their own, with discrete arrows, and whether that complex is
\emph{exact}. The answer---remarkably---is yes, and the rest of the chapter is about
why and what it buys.

\subsection{The discrete arrows are interpolation matrices}

The continuous arrows $\Grad,\Curl,\Div$ are differential operators. Their discrete
counterparts are \emph{matrices}: rectangular arrays that take the coefficient vector
of a field in one discrete space to the coefficient vector of its image in the next.
Palace builds them with a single constructor, the \emph{discrete interpolator}, whose
job is to realize one de Rham arrow at the discrete level. From the source
(\code{fespace.cpp}'s \code{BuildDiscreteInterpolator}), the constructor dispatches on
the \emph{pair} of spaces it is handed and selects the matching arrow:
\begin{center}
\begin{tabular}{@{}lll@{}}
\toprule
discrete arrow & from $\to$ to & kernel selected \\
\midrule
discrete gradient $\mathsf{G}$ & $\Hone_h\to\Hcurl_h$ & \code{GradientInterpolator} \\
discrete curl $\mathsf{C}$     & $\Hcurl_h\to\Hdiv_h$ & \code{CurlInterpolator} \\
discrete curl (2-D) & $\Hcurl_h\to\Ltwo_h$ & \code{CurlInterpolator} (native) \\
discrete divergence $\mathsf{D}$ & $\Hdiv_h\to\Ltwo_h$ & \code{DivergenceInterpolator} \\
\bottomrule
\end{tabular}
\end{center}
Each is an honest rectangular matrix: $\mathsf{G}$ has as many rows as there are
edges and as many columns as there are vertices; $\mathsf{C}$ as many rows as faces
and columns as edges; $\mathsf{D}$ as many rows as cells and columns as faces. (The
constructor refuses any pair that is not de Rham-adjacent in the correct
direction---a built-in check that the arrow even exists.) We write $\mathsf{G}$,
$\mathsf{C}$, $\mathsf{D}$ for the three matrices, reserving the continuous
$\Grad,\Curl,\Div$ for the operators they discretize.

The construction of $\mathsf{G}$ is so concrete it is worth stating outright, because
we will compute with it in \cref{wex:discrete}. The discrete gradient is, up to sign,
the \emph{incidence matrix} of the mesh: the matrix that records, for each edge, which
two vertices it joins.

\begin{definition}[The lowest-order discrete gradient]\label{def:discrete-grad}
Let the mesh have vertices indexed $1,\dots,n_v$ and edges $1,\dots,n_e$, with edge
$e$ oriented from vertex $\operatorname{tail}(e)$ to vertex $\operatorname{head}(e)$.
The lowest-order discrete gradient is the $n_e\times n_v$ matrix $\mathsf{G}$ with
\[
  \mathsf{G}_{e,v} =
  \begin{cases}
    +1 & v = \operatorname{head}(e),\\
    -1 & v = \operatorname{tail}(e),\\
    0  & \text{otherwise.}
  \end{cases}
\]
That is, $(\mathsf{G}\,\vu)_e = u_{\operatorname{head}(e)} - u_{\operatorname{tail}(e)}$:
the discrete gradient of a nodal field reads off, on each edge, the difference of the
nodal values at its two endpoints.
\end{definition}

This is exactly right as a discrete gradient. A nodal field $u$ has one number per
vertex; its continuous gradient, integrated along an edge, is
$\int_e\Grad u\cdot d\vec\ell = u(\operatorname{head}) - u(\operatorname{tail})$ by
the fundamental theorem of calculus---and the edge degree of freedom of $\Hcurl_h$ is
exactly that edge circulation (\cref{ch:functionspaces}). So the discrete gradient
$\mathsf{G}\vu$, read as edge circulations, holds the true edge circulations of
$\Grad u$. The matrix is not an approximation of the gradient; on these degrees of
freedom it is \emph{exact}.

\subsection{The commuting diagram}

To say the discrete arrows ``are'' the continuous ones, we need a precise statement,
and it is the conceptual core of the chapter. Each space comes with a canonical
\emph{interpolant}: an operator $\Pi$ that takes a smooth field and reads off its
degrees of freedom, producing the discrete field with those DoFs. There is one per
space: $\Pi_{\Hone}$ reads nodal values, $\Pi_{\Hcurl}$ reads edge circulations,
$\Pi_{\Hdiv}$ reads face fluxes, $\Pi_{\Ltwo}$ reads cell averages. The claim is that
\emph{interpolating then differentiating discretely gives the same result as
differentiating then interpolating}:
\[
  \mathsf{G}\,\Pi_{\Hone} = \Pi_{\Hcurl}\,\Grad, \qquad
  \mathsf{C}\,\Pi_{\Hcurl} = \Pi_{\Hdiv}\,\Curl, \qquad
  \mathsf{D}\,\Pi_{\Hdiv} = \Pi_{\Ltwo}\,\Div.
\]
This is the \emph{commuting-diagram property}, drawn in \cref{fig:commute}. It says
the square formed by a continuous arrow on top, a discrete arrow on the bottom, and
interpolation arrows down the sides \emph{commutes}: both ways around the square give
the same discrete field.

\begin{figure}[t]
\centering
\begin{tikzpicture}[scale=1.0, node distance=20mm and 34mm]
  % continuous row
  \node[space, fill=simBgGray, draw=simGray, minimum width=16mm] (Ha) {$\Hone$};
  \node[space, fill=simBgGray, draw=simGray, minimum width=16mm, right=of Ha] (Hb) {$\Hcurl$};
  % discrete row
  \node[space, minimum width=16mm, below=of Ha] (ha) {$\Hone_h$};
  \node[space, minimum width=16mm, below=of Hb] (hb) {$\Hcurl_h$};
  % top arrow (continuous)
  \draw[flow] (Ha) -- node[above,font=\small]{$\Grad$} (Hb);
  % bottom arrow (discrete)
  \draw[depedge] (ha) -- node[below,font=\small]{$\mathsf{G}$} (hb);
  % side arrows (interpolation, downward)
  \draw[flow, simViolet] (Ha) -- node[left,font=\scriptsize]{$\Pi_{\Hone}$} (ha);
  \draw[flow, simViolet] (Hb) -- node[right,font=\scriptsize]{$\Pi_{\Hcurl}$} (hb);
  % commuting marker
  \node[font=\Large, simGreen] at ($(Ha)!0.5!(hb)$) {$\circlearrowleft$};
  % caption-of-square
  \node[font=\scriptsize\itshape, simGray, right=2mm of hb, text width=34mm, align=left]
    {$\mathsf{G}\,\Pi_{\Hone}=\Pi_{\Hcurl}\,\Grad$\\[2pt]both paths agree};
\end{tikzpicture}
\caption[The commuting diagram]{The commuting-diagram property, shown for the
gradient square (the curl and divergence squares are identical in form). Start at a
smooth scalar field in $\Hone$ (top left). Going \emph{right then down}: differentiate
continuously with $\Grad$, then interpolate the result into the edge space with
$\Pi_{\Hcurl}$. Going \emph{down then right}: interpolate the scalar into nodal DoFs
with $\Pi_{\Hone}$, then apply the discrete gradient matrix $\mathsf{G}$. The two
paths land on the \emph{same} discrete edge field---the square commits
($\circlearrowleft$). This single property is what makes $\Hone_h,\Hcurl_h,\Hdiv_h,
\Ltwo_h$ a faithful discrete copy of the continuous complex.}
\label{fig:commute}
\end{figure}

Why does the square commute? For the gradient case it is the fundamental theorem of
calculus again. Interpolate a smooth $u$ into nodal DoFs (read its vertex values),
then apply $\mathsf{G}$: by \cref{def:discrete-grad} the result on edge $e$ is
$u(\operatorname{head}(e))-u(\operatorname{tail}(e))$. Going the other way, take the
continuous gradient $\Grad u$, then interpolate into edge DoFs (read its
circulations): the circulation on edge $e$ is
$\int_e\Grad u\cdot d\vec\ell=u(\operatorname{head}(e))-u(\operatorname{tail}(e))$.
Identical. The curl and divergence squares commute for the same reason, with Stokes'
theorem and the divergence theorem playing the role the fundamental theorem of
calculus plays here---each relates the differential operator's DoF (face flux, cell
average) to a boundary integral of the field's lower DoFs.

\subsection{Discrete exactness}
\label{sec:discrete-exact}

The commuting diagram has an immediate and decisive corollary: the discrete arrows
compose to zero, exactly as the continuous ones do.

\begin{theorem}[Discrete exactness]\label{thm:discrete-exact}
The discrete arrows satisfy the complex identities
\[
  \mathsf{C}\,\mathsf{G} = 0, \qquad \mathsf{D}\,\mathsf{C} = 0,
\]
and on a discretization of a simply-connected domain the discrete complex
$\Hone_h\xrightarrow{\mathsf{G}}\Hcurl_h\xrightarrow{\mathsf{C}}\Hdiv_h
\xrightarrow{\mathsf{D}}\Ltwo_h$ is exact:
$\operatorname{range}(\mathsf{G})=\ker(\mathsf{C})$ and
$\operatorname{range}(\mathsf{C})=\ker(\mathsf{D})$.
\end{theorem}

\begin{proof}[Why $\mathsf{C}\mathsf{G}=0$]
We give the argument for the first identity; the second is identical with $\Div\Curl$
in place of $\Curl\Grad$. Take any discrete nodal field $u_h\in\Hone_h$; we must show
$\mathsf{C}\mathsf{G}u_h=0$. The lowest-order discrete spaces have the property that
each discrete field is the interpolant of itself, so we may write $u_h=\Pi_{\Hone}u$
for the smooth representative $u$ that $u_h$ samples. Now chase the commuting diagram
twice:
\[
  \mathsf{C}\,\mathsf{G}\,u_h
  = \mathsf{C}\,\mathsf{G}\,\Pi_{\Hone}u
  = \mathsf{C}\,\Pi_{\Hcurl}\,\Grad u
  = \Pi_{\Hdiv}\,\Curl\,\Grad u
  = \Pi_{\Hdiv}\,\mathbf{0}
  = 0,
\]
using the gradient commuting square at the second step, the curl commuting square at
the third, and the \emph{continuous} identity $\Curl\Grad=0$ (\cref{lem:complex}) at
the fourth. The discrete identity is the continuous identity, pulled through the
commuting diagram.
\end{proof}

This is the heart of the matter. \emph{Discrete exactness is not a separate
miracle---it is the continuous exactness, inherited through the commuting diagram.}
The finite-element spaces of \cref{ch:functionspaces} were not chosen for accuracy
alone; they were chosen so that their canonical interpolants commute with
differentiation, which is exactly the condition that makes the discrete complex a
faithful copy of the continuous one. \Cref{fig:incidence} shows the discrete identity
$\mathsf{C}\mathsf{G}=0$ as a statement about the mesh's incidence structure.

\begin{figure}[t]
\centering
\begin{tikzpicture}[scale=1.0, every node/.style={font=\scriptsize}]
  % a triangle with a vertex, two edges, one face
  \coordinate (A) at (0,0);
  \coordinate (B) at (2.4,0);
  \coordinate (C) at (1.2,2.0);
  \fill[simBgGreen] (A)--(B)--(C)--cycle;
  \draw[simInk,thick] (A)--(B)--(C)--cycle;
  % vertices
  \fill[simBlue] (A) circle (2.4pt); \node[simBlue,below left] at (A) {$v_1$};
  \fill[simBlue] (B) circle (2.4pt); \node[simBlue,below right] at (B) {$v_2$};
  \fill[simBlue] (C) circle (2.4pt); \node[simBlue,above] at (C) {$v_3$};
  % oriented edges with circulation arrows
  \draw[-{Stealth[length=2mm]},simRust,thick] ($(A)!0.3!(B)$)--($(A)!0.7!(B)$);
  \node[simRust,below] at ($(A)!0.5!(B)$) {$e_1$};
  \draw[-{Stealth[length=2mm]},simRust,thick] ($(B)!0.3!(C)$)--($(B)!0.7!(C)$);
  \node[simRust,right] at ($(B)!0.5!(C)$) {$e_2$};
  \draw[-{Stealth[length=2mm]},simRust,thick] ($(C)!0.3!(A)$)--($(C)!0.7!(A)$);
  \node[simRust,left] at ($(C)!0.5!(A)$) {$e_3$};
  % face circulation marker
  \draw[simGreen,thick,-{Stealth[length=1.8mm]}] (1.2,0.75) arc (-90:170:0.32);
  \node[simGreen] at (1.2,0.62) {$f$};
  % the algebra to the right
  \node[align=left, font=\scriptsize, text width=58mm, anchor=west] at (3.4,1.0) {%
    \textbf{Discrete gradient} $\mathsf{G}$ (edge $=$ head $-$ tail):\\[2pt]
    $\mathsf{G}\vu=(u_2{-}u_1,\;u_3{-}u_2,\;u_1{-}u_3)$\\[4pt]
    \textbf{Discrete curl} $\mathsf{C}$ (face $=$ signed edge sum):\\[2pt]
    $\mathsf{C}(\mathsf{G}\vu)=(u_2{-}u_1)+(u_3{-}u_2)+(u_1{-}u_3)=0$\\[4pt]
    The loop of edge differences \emph{telescopes}: $\mathsf{C}\mathsf{G}=0$.};
\end{tikzpicture}
\caption[Discrete $\mathsf{C}\mathsf{G}=0$ telescopes]{The discrete identity
$\mathsf{C}\mathsf{G}=0$ on a single triangle. The discrete gradient places, on each
oriented edge, the difference of its endpoint nodal values (head minus tail). The
discrete curl sums the edge circulations around the face boundary, with sign. Going
around the loop $e_1,e_2,e_3$, the differences \emph{telescope}: every vertex value
appears once with $+$ and once with $-$, so the sum is identically zero. This is the
discrete avatar of ``the line integral of a gradient around a closed loop is
zero''---and the discrete avatar of $\Curl\Grad=0$.}
\label{fig:incidence}
\end{figure}

\begin{keyidea}
\textbf{Discrete exactness is why edge elements avoid spurious modes.} The
spurious-mode pathology of \cref{ch:functionspaces} had a structural cause: a
componentwise-nodal discretization of $\vE$ has the \emph{wrong} kernel for its
discrete curl, so the discrete curl--curl operator's null space is polluted with
non-physical fields that masquerade as resonances. Edge elements ($\Hcurl_h$) fix
this because they sit in an \emph{exact} discrete complex: by
\cref{thm:discrete-exact} the kernel of the discrete curl is \emph{exactly} the range
of the discrete gradient---the discrete gauge fields $\mathsf{G}\,\psi_h$, and nothing
else. The curl--curl operator's null space is then precisely the discrete gradients,
a space we can identify and handle (\cref{sec:divfree}), with no spurious extras
hiding in it. Get the discrete complex exact and the spectrum is clean; break
exactness and spurious modes return. This is the single deepest reason the four
finite-element spaces are the ones a Maxwell solver must use.
\end{keyidea}

\section{Whitney forms: the canonical basis}
\label{sec:whitney}

The lowest-order discrete complex has a canonical basis---one set of shape functions,
one per mesh entity, that realizes all four spaces at once and makes the discrete
arrows the incidence matrices of \cref{def:discrete-grad}. These are the
\emph{Whitney forms}. We met one already: the edge function
$\boldsymbol\phi_{ij}=\lambda_i\Grad\lambda_j-\lambda_j\Grad\lambda_i$ of
\cref{wex:whitney}. The full family, built from the barycentric coordinates
$\lambda_0,\dots,\lambda_d$ of a simplex, is:

\begin{itemize}[nosep]
\item \textbf{Nodal} ($\Hone_h$, on vertices): $\lambda_i$ itself---the barycentric
  hat function, $1$ at vertex $i$, $0$ at the others.
\item \textbf{Edge} ($\Hcurl_h$, on edges): $\boldsymbol\phi_{ij}=
  \lambda_i\Grad\lambda_j-\lambda_j\Grad\lambda_i$ on edge $(i,j)$---unit circulation
  on its own edge, zero on the others.
\item \textbf{Face} ($\Hdiv_h$, on faces): $\boldsymbol\phi_{ijk}=
  2\bigl(\lambda_i(\Grad\lambda_j\times\Grad\lambda_k)
  +\lambda_j(\Grad\lambda_k\times\Grad\lambda_i)
  +\lambda_k(\Grad\lambda_i\times\Grad\lambda_j)\bigr)$ on face $(i,j,k)$---unit flux
  through its own face.
\item \textbf{Cell} ($\Ltwo_h$, on cells): the constant (or $\det$-scaled constant)
  on each cell---unit cell average.
\end{itemize}

\begin{figure}[t]
\centering
\begin{tikzpicture}[scale=1.05, every node/.style={font=\scriptsize}]
% three triangles showing nodal / edge / face Whitney forms
% --- nodal ---
\begin{scope}
  \node[font=\footnotesize\bfseries] at (1.0,2.1) {nodal $\lambda_i$};
  \coordinate (a) at (0,0); \coordinate (b) at (2.0,0); \coordinate (c) at (1.0,1.7);
  \draw[simInk,thick] (a)--(b)--(c)--cycle;
  \fill[simBlue] (a) circle (3pt);
  \fill[simBlue!35] (b) circle (1.6pt); \fill[simBlue!35] (c) circle (1.6pt);
  \node[simGray] at (1.0,-0.45) {one per vertex};
\end{scope}
% --- edge ---
\begin{scope}[shift={(3.6,0)}]
  \node[font=\footnotesize\bfseries] at (1.0,2.1) {edge $\boldsymbol\phi_{ij}$};
  \coordinate (a) at (0,0); \coordinate (b) at (2.0,0); \coordinate (c) at (1.0,1.7);
  \draw[simInk,thick] (a)--(b)--(c)--cycle;
  % tangential circulation along the bottom edge
  \foreach \t in {0.2,0.4,0.6,0.8}{
    \draw[-{Stealth[length=1.5mm]},simRust,thick]
      ($(a)!\t!(b)+(0,0.12)$)--+(0.28,0);}
  \node[simGray] at (1.0,-0.45) {one per edge};
\end{scope}
% --- face ---
\begin{scope}[shift={(7.2,0)}]
  \node[font=\footnotesize\bfseries] at (1.0,2.1) {face $\boldsymbol\phi_{ijk}$};
  \coordinate (a) at (0,0); \coordinate (b) at (2.0,0); \coordinate (c) at (1.0,1.7);
  \fill[simGreen!16] (a)--(b)--(c)--cycle;
  \draw[simInk,thick] (a)--(b)--(c)--cycle;
  % flux out of the face (dot = out of page)
  \node[circle,draw=simGreen,fill=white,inner sep=1.1pt] at (1.0,0.55) {};
  \fill[simGreen] (1.0,0.55) circle (1.0pt);
  \node[simGreen] at (1.45,0.55) {$\odot$};
  \node[simGray] at (1.0,-0.45) {one per face};
\end{scope}
\end{tikzpicture}
\caption[Whitney forms on a triangle]{The lowest-order Whitney basis, the canonical
realization of the discrete de Rham complex. \emph{Left:} the nodal form $\lambda_i$,
one per vertex, peaking at its vertex (the $\Hone_h$ basis). \emph{Middle:} the edge
form $\boldsymbol\phi_{ij}$, one per edge, carrying unit tangential circulation along
its own edge (the $\Hcurl_h$ basis of \cref{wex:whitney}). \emph{Right:} the face form
$\boldsymbol\phi_{ijk}$, one per face, carrying unit normal flux through its own face
(the $\Hdiv_h$ basis, drawn here as flux out of the page, $\odot$). Together they make
the discrete arrows the incidence matrices of \cref{def:discrete-grad}: the gradient
of a nodal hat is a sum of edge forms, and so on up the complex.}
\label{fig:whitney}
\end{figure}

The decisive algebraic fact---the one that ties the basis to the incidence matrices---
is how the differential operators act on the Whitney forms. Applying the gradient to a
nodal hat $\lambda_i$ produces a field that, expanded in the edge forms
$\boldsymbol\phi_{ij}$, has coefficients $\pm 1$ on the edges touching vertex $i$ and
$0$ elsewhere---exactly the column of the incidence matrix $\mathsf{G}$ for vertex
$i$. The Whitney forms are the basis in which the discrete gradient, curl, and
divergence \emph{are} the signed incidence matrices of the mesh. This is why
\cref{def:discrete-grad} could state $\mathsf{G}$ combinatorially: in the Whitney
basis, the analysis is purely topological.

\begin{remark}[For the curious: cohomology and holes]
\label{rmk:cohomology}
We claimed exactness for \emph{simply-connected} domains and warned of a gap
otherwise. That gap is not a defect---it is information, and it is the gateway to a
beautiful corner of mathematics. On a domain with holes, the quotient
$\ker(\Curl)/\operatorname{range}(\Grad)$---the curl-free fields that are \emph{not}
gradients, modulo those that are---is a finite-dimensional vector space whose
dimension \emph{counts the holes}. Its elements are the \emph{harmonic} fields: those
in the kernel of both the differential operator above and the (adjoint) one below.
For a solid torus this space is one-dimensional; for a domain with $g$ independent
tunnels it is $g$-dimensional. This is \emph{de Rham cohomology}, and its punchline is
startling: a purely analytic object (which curl-free fields fail to be gradients) is
counting a purely topological one (how many holes the domain has). The discrete
complex inherits this too---a correctly-constructed Whitney complex on a meshed torus
has exactly one discrete harmonic field, and a solver can find it. Palace's domains
are usually simply-connected, so the harmonic space is trivial and the complex is
exact; but the machinery is there, and the fact that finite elements can count the
holes of a domain is one of the quiet triumphs of the theory. The canonical reference
is Arnold, Falk, and Winther's finite-element exterior calculus~\citep{arnold2006feec};
we say no more here.
\end{remark}

\section{Worked examples}

\subsection{The continuous identities, component by component}

We first watch \cref{lem:complex} happen on concrete polynomial fields, so the
mixed-partial cancellation is not an abstraction but an arithmetic fact you can check.

\begin{workedexample}{Verifying $\Curl\Grad\phi=0$ and $\Div\Curl\vF=0$}{component}
\textbf{Part 1: $\Curl\Grad\phi=0$.} Take the scalar field
$\phi(x,y,z)=x^2 y + y z^2$. Its gradient is
\[
  \Grad\phi = \bigl(\,2xy,\;\; x^2 + z^2,\;\; 2yz\,\bigr).
\]
Now take the curl of this vector field. The $x$-component is
\[
  (\Curl\Grad\phi)_x
  = \partial_y(2yz) - \partial_z(x^2+z^2)
  = 2z - 2z = 0.
\]
The $y$-component is
\[
  (\Curl\Grad\phi)_y
  = \partial_z(2xy) - \partial_x(2yz)
  = 0 - 0 = 0.
\]
The $z$-component is
\[
  (\Curl\Grad\phi)_z
  = \partial_x(x^2+z^2) - \partial_y(2xy)
  = 2x - 2x = 0.
\]
So $\Curl\Grad\phi=\mathbf 0$. Notice the mechanism in the $x$-component: the term
$2z$ came from $\partial_y\partial_z\phi$ and the term $-2z$ from
$-\partial_z\partial_y\phi$; they cancel because mixed partials commute. The other
components cancel the same way.

\textbf{Part 2: $\Div\Curl\vF=0$.} Take the vector field
$\vF(x,y,z)=\bigl(\,xy z,\;\; y^2 x,\;\; z x^2\,\bigr)$. Its curl is
\[
  \Curl\vF =
  \begin{pmatrix}
    \partial_y(z x^2) - \partial_z(y^2 x)\\[2pt]
    \partial_z(x y z) - \partial_x(z x^2)\\[2pt]
    \partial_x(y^2 x) - \partial_y(x y z)
  \end{pmatrix}
  =
  \begin{pmatrix}
    0 - 0\\
    x y - 2 z x\\
    y^2 - x z
  \end{pmatrix}
  =
  \begin{pmatrix}
    0\\ xy - 2zx\\ y^2 - xz
  \end{pmatrix}.
\]
Now its divergence:
\[
  \Div\Curl\vF
  = \partial_x(0) + \partial_y(xy - 2zx) + \partial_z(y^2 - xz)
  = 0 + x + (-x)
  = 0.
\]
Again the surviving terms cancel in a mixed-partial pair: $\partial_y(xy)=x$ from the
second component against $\partial_z(-xz)=-x$ from the third. The identity holds not
by luck but by the structure of the curl, which arranges every term to meet its
negative.
\end{workedexample}

\subsection{Discrete $\mathsf{C}\mathsf{G}=0$ on a small mesh}

Now the discrete version. We build the lowest-order discrete gradient on a tiny
two-triangle mesh, apply it to a concrete nodal field, and verify by hand that the
discrete curl annihilates the result---a single instance of
\cref{thm:discrete-exact}.

\begin{workedexample}{The discrete gradient on a square of two triangles}{discrete}
\textbf{The mesh.} Take a unit square split into two triangles by its diagonal, with
vertices
\[
  v_1=(0,0),\quad v_2=(1,0),\quad v_3=(1,1),\quad v_4=(0,1),
\]
and five edges---the four sides and the diagonal---oriented as
\[
  e_1: v_1\!\to\!v_2,\;\;
  e_2: v_2\!\to\!v_3,\;\;
  e_3: v_4\!\to\!v_3,\;\;
  e_4: v_1\!\to\!v_4,\;\;
  e_5: v_1\!\to\!v_3 \ (\text{diagonal}).
\]
The two triangular faces are
$f_1=(v_1,v_2,v_3)$ and $f_2=(v_1,v_3,v_4)$.

\textbf{The discrete gradient $\mathsf{G}$.} By \cref{def:discrete-grad}, each row of
the $5\times 4$ matrix $\mathsf{G}$ is $+1$ at the edge's head, $-1$ at its tail:
\[
  \mathsf{G} =
  \begin{array}{c@{\,}c}
  & \begin{matrix} v_1 & v_2 & v_3 & v_4 \end{matrix}\\
  \begin{matrix} e_1\\ e_2\\ e_3\\ e_4\\ e_5 \end{matrix} &
  \begin{pmatrix*}[r]
    -1 & 1 & 0 & 0\\
     0 & -1 & 1 & 0\\
     0 & 0 & 1 & -1\\
    -1 & 0 & 0 & 1\\
    -1 & 0 & 1 & 0
  \end{pmatrix*}
  \end{array}.
\]
\textbf{A nodal field and its discrete gradient.} Take the nodal field
$\vu=(u_1,u_2,u_3,u_4)=(0,1,3,2)$ (one value per vertex). Then
\[
  \mathsf{G}\vu =
  \begin{pmatrix*}[r]
    u_2-u_1\\ u_3-u_2\\ u_3-u_4\\ u_4-u_1\\ u_3-u_1
  \end{pmatrix*}
  =
  \begin{pmatrix*}[r]
    1-0\\ 3-1\\ 3-2\\ 2-0\\ 3-0
  \end{pmatrix*}
  =
  \begin{pmatrix*}[r] 1\\ 2\\ 1\\ 2\\ 3 \end{pmatrix*}.
\]
These five numbers are the edge circulations of the discrete gradient field
$\Grad u_h$: one per edge.

\textbf{The discrete curl, face by face, and $\mathsf{C}\mathsf{G}=0$.} The discrete
curl on a face is the signed sum of the circulations of its boundary edges, traversed
consistently with the face orientation. Walk the boundary of $f_1=(v_1,v_2,v_3)$:
$v_1\!\to\!v_2$ along $+e_1$, $v_2\!\to\!v_3$ along $+e_2$, and $v_3\!\to\!v_1$ along
$-e_5$ (the diagonal, traversed against its orientation). So
\[
  (\mathsf{C}\mathsf{G}\vu)_{f_1}
  = (+1)(\mathsf{G}\vu)_{e_1} + (+1)(\mathsf{G}\vu)_{e_2} + (-1)(\mathsf{G}\vu)_{e_5}
  = 1 + 2 - 3 = 0.
\]
Now the other face $f_2=(v_1,v_3,v_4)$: $v_1\!\to\!v_3$ along $+e_5$, $v_3\!\to\!v_4$
along $-e_3$, and $v_4\!\to\!v_1$ along $-e_4$. So
\[
  (\mathsf{C}\mathsf{G}\vu)_{f_2}
  = (+1)(\mathsf{G}\vu)_{e_5} + (-1)(\mathsf{G}\vu)_{e_3} + (-1)(\mathsf{G}\vu)_{e_4}
  = 3 - 1 - 2 = 0.
\]
Both faces give zero: $\mathsf{C}\mathsf{G}\vu=\mathbf 0$. And the reason is exactly
the telescoping of \cref{fig:incidence}: on $f_1$ the circulations are
$(u_2-u_1)+(u_3-u_2)-(u_3-u_1)$, in which every nodal value cancels its own appearance.
The result holds for \emph{every} nodal field $\vu$, not just our chosen one---that
is the content of $\mathsf{C}\mathsf{G}=0$. The discrete gradient of any nodal function
has zero discrete curl, edge differences around any face loop always telescoping to
nothing.
\end{workedexample}

\section{Consequences for the solver}
\label{sec:divfree}

The de Rham structure is not theory for its own sake; the solver chapters lean on it
constantly. We close by collecting the three consequences that recur, each a direct
reading of discrete exactness.

\paragraph{The gradient null space of curl--curl.} Every magnetic and wave analysis
assembles the curl--curl operator $\mathsf{K}\sim\Curl(\nu\Curl\,\cdot)$ on $\Hcurl_h$.
By discrete exactness, $\mathsf{C}\mathsf{G}=0$, so this operator annihilates every
discrete gradient: $\mathsf{K}\,\mathsf{G}\psi_h=0$ for any nodal $\psi_h$. The
operator is therefore \emph{singular}, with null space exactly the discrete gradients
$\operatorname{range}(\mathsf{G})$---the discrete gauge fields of \cref{sec:potentials}.
This is not a bug to be regularized away; it is the discrete image of gauge freedom,
and discrete exactness tells us precisely what it is.

\paragraph{Projecting onto the divergence-free subspace.} Because the curl--curl
operator has this gradient null space, an eigenmode solve will, left alone, return
spurious ``zero-frequency'' modes living in $\operatorname{range}(\mathsf{G})$,
interleaved with the physical resonances. The fix is to confine the solver to the
\emph{complement}: the divergence-free subspace, those Nedelec fields $\vy$ with
$\mathsf{G}^\top\mathsf{M}\vy=0$ (no discrete-gradient component in the $\mathsf{M}$-inner
product). Palace's divergence-free projector
\[
  P = I - \mathsf{G}\,(\mathsf{G}^\top\mathsf{M}\mathsf{G})^{-1}\mathsf{G}^\top\mathsf{M}
\]
removes the gradient part of any field, mapping it into the divergence-free subspace;
the eigensolver applies $P$ to each candidate vector to keep it physical. Every
ingredient of $P$ is a de Rham object: $\mathsf{G}$ is the discrete gradient built by
the interpolator of \cref{sec:discrete-arrows}, and the subspace it projects onto is
the kernel of the discrete divergence the complex defines. We develop the projector and
its place in the eigenmode pipeline in \cref{ch:eigenmodes}.

\paragraph{The Helmholtz decomposition.} Underlying both is the discrete
\emph{Helmholtz decomposition}: every Nedelec field splits uniquely as a
divergence-free part plus a gradient part,
\[
  \vy = \underbrace{\vy_{\mathrm{df}}}_{\text{divergence-free}}
        \;+\; \underbrace{\mathsf{G}\psi}_{\text{gradient}},
\]
and the projector $P$ is exactly the operator that extracts the first summand. This is
the discrete shadow of the continuous decomposition (\cref{ch:bc}), and it is a direct
consequence of discrete exactness: the two summands live in $\ker(\mathsf{C})$'s
complement and in $\operatorname{range}(\mathsf{G})=\ker(\mathsf{C})$ respectively, the
equality being \cref{thm:discrete-exact}. The same decomposition is what makes the
magnetostatic and electrostatic analyses well-posed, and it is the structural reason the
boundary-condition chapter can speak of ``the divergence-free field'' as a definite
object.

\begin{keyidea}
Three solver facts, one structural source. The curl--curl operator is singular
\emph{because} $\mathsf{C}\mathsf{G}=0$; its null space is the discrete gradients
\emph{because} the discrete complex is exact; and the eigensolver projects onto the
divergence-free subspace \emph{because} that is the exactness-supplied complement of
the gradient null space. Discrete exactness is not a property one checks and files
away---it is load-bearing in every magnetic, wave, and eigenmode solve in the book.
\end{keyidea}

\summarysection
The four finite-element spaces of \cref{ch:functionspaces} assemble into the
\emph{de Rham complex}
$\Hone\xrightarrow{\Grad}\Hcurl\xrightarrow{\Curl}\Hdiv\xrightarrow{\Div}\Ltwo$:
each differential operator maps one space exactly into the next. The complex
\emph{property} is the pair of identities $\Curl\Grad=0$ and $\Div\Curl=0$, both
corollaries of the commuting of mixed partial derivatives; \emph{exactness} is the
stronger statement that the range of each arrow equals the kernel of the next, true on
simply-connected domains. Exactness is the mathematical content of the electromagnetic
potentials: $\Curl\vE=0$ forces $\vE=-\Grad V$, and $\Div\vB=0$ forces $\vB=\Curl\vA$,
the potential always living one rung up the complex. The discrete finite-element spaces
inherit the same structure: their canonical interpolants \emph{commute with
differentiation} (the commuting-diagram property), and through that diagram the discrete
arrows---the incidence-style matrices $\mathsf{G},\mathsf{C},\mathsf{D}$ built by
Palace's discrete interpolator---satisfy $\mathsf{C}\mathsf{G}=0$ and
$\mathsf{D}\mathsf{C}=0$ and form an \emph{exact discrete complex}. The Whitney forms are
the canonical lowest-order basis realizing it, in which the discrete arrows are the
signed incidence matrices of the mesh. Discrete exactness is the deep reason edge
elements avoid spurious modes: it makes the curl--curl operator's null space exactly the
discrete gradients, no spurious extras, so the eigensolver can project onto the
divergence-free complement and recover a clean spectrum. Every divergence-free
projection, every gauge null space, and the Helmholtz decomposition of \cref{ch:bc} is a
reading of this one structure.

\furtherreading
\label{sec:furtherreading}
The unifying modern framework is \emph{finite-element exterior calculus}, which casts
the four spaces as discrete differential forms and the complex as the central object;
Arnold, Falk, and Winther~\citep{arnold2006feec} is the canonical and remarkably readable
account, and is the natural sequel to this chapter for the cohomology remark
(\cref{rmk:cohomology}). For the electromagnetic specialization---$\Hcurl$, $\Hdiv$, the
N\'ed\'elec and Raviart--Thomas elements, the commuting interpolants, and the spurious-mode
analysis worked through with full rigour---Monk~\citep{monk2003} is the standard reference.
The discrete de Rham complex and the role of Whitney forms in computational
electromagnetics were brought to prominence by Hiptmair, whose work on the canonical
basis and the commuting-diagram property~\citep{hiptmair1998} is the bridge between the
abstract complex and the matrices a solver actually builds. Readers wanting the continuous
exactness theorem (\cref{thm:exact}, the Poincar\'e lemma) and the topology of the harmonic
spaces will find both in any of the three.

\exercisessection
\begin{enumerate}[nosep]
\item \textbf{Type-checking the arrows.} For the 2-D complex
  $\Hone\xrightarrow{\Grad}\Hcurl\xrightarrow{\operatorname{curl}}\Ltwo$, write out the
  scalar curl $\operatorname{curl}(u_1,u_2)=\partial_x u_2-\partial_y u_1$ and verify
  directly that $\operatorname{curl}(\Grad\phi)=0$ for any scalar $\phi$. Which single
  mixed-partial cancellation does it rest on?
\item \textbf{Another concrete instance.} Repeat \cref{wex:component} for
  $\phi=e^{x}\sin(yz)$: compute $\Grad\phi$ and confirm $\Curl\Grad\phi=\mathbf 0$
  component by component. (You will need the product and chain rules, but the
  cancellation is the same mixed-partial symmetry.)
\item \textbf{Range versus kernel.} On the annulus $\{1<x^2+y^2<2\}$ in the plane, the
  field $\vF=(-y,x)/(x^2+y^2)$ is curl-free but is \emph{not} the gradient of any
  single-valued scalar. Compute $\operatorname{curl}\vF$ to confirm it is curl-free, then
  explain in one sentence why this does not contradict \cref{thm:exact}. (Which
  hypothesis of the theorem does the annulus violate?)
\item \textbf{Potentials are exactness.} State, in the language of ranges and kernels,
  the precise de Rham fact that licenses each of: (a) $\vE=-\Grad V$ in electrostatics;
  (b) $\vB=\Curl\vA$ in magnetostatics. For each, name the rung of the complex at which
  the exactness is read.
\item \textbf{Gauge freedom.} Show that if $\vB=\Curl\vA$ then also $\vB=\Curl\vA'$ for
  $\vA'=\vA+\Grad\psi$, any scalar $\psi$. Which de Rham identity makes this work, and how
  does it explain why the discrete curl--curl operator is singular?
\item \textbf{Discrete gradient on a hexahedral face.} Build the lowest-order discrete
  gradient $\mathsf{G}$ for a single square cell with vertices
  $v_1,v_2,v_3,v_4$ (counterclockwise) and four boundary edges oriented
  $v_1\!\to\!v_2\!\to\!v_3\!\to\!v_4\!\to\!v_1$. Take the nodal field
  $\vu=(2,5,4,1)$, compute the four edge circulations $\mathsf{G}\vu$, and verify that
  the discrete curl on the single face (the signed sum around the boundary) is zero.
\item \textbf{Telescoping is the mechanism.} Prove the general claim behind
  \cref{wex:discrete}: for \emph{any} oriented cycle of edges
  $e_1,\dots,e_k$ forming a closed loop and any nodal field $\vu$, the signed sum of
  edge differences $\sum_j \pm(u_{\mathrm{head}(e_j)}-u_{\mathrm{tail}(e_j)})$ around the
  loop is zero. (This is $\mathsf{C}\mathsf{G}=0$ for one face; the sign on $e_j$ is its
  orientation relative to the loop traversal.)
\item \textbf{The commuting square in symbols.} Using the commuting-diagram property
  $\mathsf{G}\,\Pi_{\Hone}=\Pi_{\Hcurl}\,\Grad$ and the analogous curl square, reproduce
  the diagram-chase proof that $\mathsf{C}\mathsf{G}=0$, and identify which single
  \emph{continuous} fact (\cref{lem:complex}) the discrete identity inherits.
\item \textbf{Why not nodal vectors, revisited.} In one paragraph, connect the
  spurious-mode pitfall of \cref{ch:functionspaces} to discrete exactness: explain why a
  componentwise-nodal discretization of $\vE$ fails to sit in an exact discrete complex,
  and what goes wrong with the kernel of its discrete curl as a result.
\item \textbf{Cohomology counts holes (open-ended).} Read \cref{rmk:cohomology} again. On
  a solid torus, how many independent discrete harmonic fields does the Whitney complex
  carry, and what feature of the geometry does that number measure? (No computation
  required---state the count and its topological meaning.)
\end{enumerate}
